\begin{abstract}
Robotic garment folding requires policies that generalize across unseen garment styles, yet all three baselines provided by the LeHome Challenge at ICRA 2026 rely on RGB-based visual encoders that encode color and texture rather than geometry. We argue that this is the root cause of poor cross-style generalization, since two shirts of different styles share almost no visual similarity in appearance despite having consistent underlying geometric structure. We extend 3D Diffusion Policy (DP3) with Normalized Object Coordinate Space (NOCS) normalization, replacing RGB inputs with depth-derived point clouds and mapping garment geometry into a shared canonical space. Our approach is motivated by two complementary findings in prior work: DP3 demonstrated categorical generalization advantages over RGB-based policies on rigid objects via point cloud representations, and UniFolding demonstrated that NOCS normalization enables cross-garment generalization on folding tasks. Whether these advantages transfer to deformable garment folding with a continuous joint trajectory action space is an open question that our project directly addresses. We evaluate through a structured ablation study on the LeHome simulation benchmark, measuring success rate on held-out unseen garment styles across four categories.
\end{abstract}
